The rapid growth of open-source software (OSS) has transformed the way software is developed, enabling globally distributed collaboration and widespread access to code and innovation. Platforms like GitHub have become central to this ecosystem, hosting millions of projects and developers who contribute in diverse ways—from reporting issues and engaging in discussions to submitting code changes. However, despite the openness of OSS communities, not all contributors transition into core team members. Many projects struggle to balance inclusivity with maintaining high-quality contributions, making it difficult to determine which types of developer interactions actually lead to team acceptance. This challenge highlights a critical gap in understanding how contributors evolve from peripheral participants to trusted collaborators within OSS teams.

To address this gap, the paper investigates the patterns of developer activity on GitHub that distinguish those who successfully join project teams from those who remain outsiders. By analyzing large-scale interaction data—including pull requests, comments, and other forms of engagement—the study aims to identify which behaviors are most predictive of joining outcomes. Rather than focusing solely on code contributions, the work considers a broader spectrum of social and technical interactions, emphasizing the role of communication and visibility in collaborative environments. Ultimately, the study seeks to provide insights that can inform OSS platform design, community practices, and developer education, helping both contributors and teams better navigate the process of onboarding and collaboration.

\section{Motivation}
The motivation for this paper arises from the increasing importance of open-source software (OSS) in modern software development and the growing reliance of society on OSS-based systems. The proliferation of the Internet and OSS philosophies has enabled developers across the globe to collaboratively build and maintain large-scale systems, such as the Linux kernel and Mozilla Firefox \cite{1_Nakakoji2003,2_OSDDefinition,3_Fitzgerald2006}. While this openness lowers barriers to entry and encourages widespread participation, it also introduces challenges in sustaining these projects. Many OSS teams struggle to attract and retain sufficient contributors, which can lead to serious consequences for software quality and security. A notable example is the Heartbleed vulnerability in OpenSSL, where a small, overstretched team failed to detect a critical bug affecting a large portion of the web \cite{4_Durumeric2014}. This illustrates that simply having open access to contributions is not enough; projects must also ensure that contributors are capable and trustworthy, creating a tension between inclusivity and quality control \cite{5_Bird2007}.

This tension motivates a deeper investigation into how developers transition from being external contributors to becoming trusted members of OSS teams. Although collaborative platforms like GitHub provide mechanisms for interaction—such as pull requests, issue discussions, and social signals—they also create a complex ecosystem where it is difficult to identify which behaviors actually lead to team acceptance. With millions of users and projects, GitHub serves as a rich environment where diverse forms of developer activity can be observed at scale. However, despite the abundance of data, there remains a lack of clarity on what distinguishes successful contributors from those who remain on the periphery. This gap in understanding limits both developers, who seek to integrate into projects, and teams, who aim to identify and recruit valuable contributors effectively.

Furthermore, prior research has explored OSS communities through structural and social lenses, such as core-periphery models and joining scripts, but often lacks quantitative analysis of the specific contribution behaviors that drive team acceptance \cite{1_Nakakoji2003}. Existing studies frequently emphasize social connections or initial interactions rather than the sustained patterns of activity that characterize successful contributors. As a result, there is a need for a more comprehensive and data-driven understanding of how different types of contributions—ranging from passive attention to active code submissions and discussions—impact the likelihood of joining a team. This paper is motivated by the opportunity to leverage large-scale GitHub data to systematically analyze these behaviors and uncover actionable insights.

Ultimately, the motivation extends beyond academic curiosity to practical implications for the OSS ecosystem. By identifying which forms of interaction most strongly influence team acceptance, the study aims to inform the design of collaborative platforms, guide educators in preparing developers for OSS participation, and help contributors adopt more effective engagement strategies. Understanding these dynamics is essential for improving the sustainability and resilience of OSS projects, which increasingly underpin critical infrastructure and software systems worldwide.

\subsection{Related Work}
\subsubsection{The Structures of OSS Teams}
Research on OSS communities commonly models their structure using a core periphery or “onion” model, where developers are organized into layers based on their level of involvement and responsibility. At the outermost layers are passive users and occasional contributors, while the innermost layers consist of core developers and project leaders who actively maintain the system \cite{1_Nakakoji2003}. Studies have shown that this structure is critical for the health of OSS projects, as it balances participation and responsibility across different roles \cite{6_Mockus2002,7_Dinh2005}. Importantly, contributions do not just modify the codebase but also influence the social structure of the community by redefining contributors' roles and positions within it \cite{1_Nakakoji2003}. However, some researchers argue that this model does not fully capture the dynamic nature of OSS communities, where developers frequently move between layers rather than remaining fixed in one position \cite{8_Ducheneaut2005}.

Further work highlights that membership in OSS teams is not uniformly defined and varies across projects. Some perspectives consider anyone who contributes or provides meaningful input as part of the team, while others emphasize formal roles such as commit access or “committer” status as indicators of true membership \cite{11_Vasilescu2015,9_Shah2006}. Additionally, studies suggest that core developers are typically those who contribute significant and frequent code changes, although formal designations do not always align perfectly with perceived roles \cite{12_Souza2005,13_Crowston2017,14_Joblin2016}. These variations underscore the complexity of defining team boundaries in OSS and motivate the need for standardized, observable definitions of membership—such as role assignments on platforms like GitHub—to enable large-scale empirical analysis.

\subsubsection{The Theories of OSS Joining}
A key line of research investigates how developers transition from peripheral contributors to core team members. One influential concept is the “joining script,” which describes a general sequence of activities that newcomers follow to gain acceptance into a project \cite{15_vonKrogh2003}. This script is not rigid but reflects common behavioral patterns, such as starting with bug reporting or minor contributions and gradually increasing involvement. Empirical studies support this idea, showing that successful joiners often engage in specific activities—like fixing bugs or contributing to discussions—while non-joiners may exhibit more limited or less targeted participation \cite{16_SinhaMS11}. These findings suggest that joining is shaped by both the type and progression of contributions rather than isolated actions.

Other research expands this perspective by examining factors beyond direct contributions, including social connections, prior experience, and visibility within the community. For example, having existing ties to core developers or prior contributions to other OSS projects significantly increases the likelihood of joining \cite{16_SinhaMS11,17_CasalnuovoVDF15}. Additionally, early engagement behaviors, such as watching projects, can indicate future involvement \cite{18_SheoranBKDE14}. However, not all contributors follow the same trajectory; differences exist between volunteer contributors and those employed by organizations to work on OSS, leading to distinct joining patterns \cite{19_HerraizRARG06}. Despite these insights, much of the prior work focuses on qualitative analysis or specific case studies, leaving a gap in quantitatively understanding how different types of contributions—especially their frequency and consistency—affect joining outcomes.

\subsubsection{Using GitHub for OSS Research}
GitHub has become a central platform for studying OSS due to its transparency and the richness of interaction data it provides. Researchers have emphasized the concept of “social translucence,” where visible traces of user activity—such as commits, pull requests, and comments—allow others to infer a developer's skills, interests, and level of engagement \cite{20_DabbishSTH12}. These activity traces play a critical role in shaping how contributors are perceived and evaluated within the community. Additional studies show that users form impressions based on both technical contributions and social signals, which influence decisions such as accepting code changes or inviting contributors to collaborate more closely \cite{21_MarlowDH13}.

Beyond individual evaluation, GitHub also enables researchers to study how social and technical dynamics influence broader community behavior. For instance, prior work has examined how norms, such as testing practices, spread within projects and how social interactions facilitate or hinder this process \cite{22_PhamSLFS13}. However, much of this research relies on qualitative methods, such as interviews or case studies, to understand how GitHub's interface shapes collaboration. In contrast, there is a need for large-scale quantitative analyses that connect these observable activities to concrete outcomes, such as team formation and developer onboarding. This gap motivates the current study's use of GitHub as a unified platform to systematically analyze contribution patterns across thousands of projects.

\subsection{Research Questions}
As the influence of OSS projects grows and more consumers depend on their outcomes, the necessity of keeping these projects well-maintained becomes all the more critical. But we do not know of,nor is there likely to be an exact and universal process to join development teams; the requirements for joining a software team varies between project types and team culture. The authors pose three research questions
\begin{enumerate}[label=\textbf{RQ\arabic*:}, leftmargin=*]
    \item What types and qualities of project interactions distinguish between developers who join an OSS team and developers who remain outside?
    \item What activities or qualities of a developer independent from a specific team influence whether the developer joins that team?
    \item What activities or qualities of a team influence whether the team accepts more people in?
\end{enumerate}

